\chapter{Harnessing Layer 1: Hard Representations and Logistics}

\section{Control Theory and The Physical API}
A "Hard Representation" is a strictly verified physical construct that reliably performs logic, structure, or exergy transfer within Layer 1. To bridge the gap from a passive digital twin to an active, steerable sovereign node, the Oasis engine must define a rigorous bidirectional interface (the "Physical API") between the computational layers (L3-L6) and the physical substrate (L1). 

This interface is governed by classic Control Theory. The computational orchestrator cannot act upon the world instantaneously; it relies on Layer 2 telemetry (sensors and actuators). The accuracy of the digital twin's state representation is bounded by the Nyquist-Shannon sampling theorem:
\begin{equation}
f_s > 2B
\end{equation}
where the sampling rate $f_s$ of the IoT sensors (e.g., thermocouples, LiDAR, current transformers) must be more than twice the maximum frequency $B$ of the physical phenomenon being monitored. If the sampling rate is insufficient, the engine will experience aliasing, leading to catastrophic mismanagement of critical workflows like Scenario Theta (Foundry Kilns) or Scenario Upsilon (Microgrid Power Shedding) \cite{shannon1949communication}.

To actively steer the physical node toward a desired intent state, the orchestrator utilizes Proportional-Integral-Derivative (PID) control algorithms:
\begin{equation}
u(t) = K_p e(t) + K_i \int_{0}^{t} e(\tau) d\tau + K_d \frac{de(t)}{dt}
\end{equation}
where $e(t)$ is the error between the desired intent (e.g., maintaining a $20^\circ$C ambient temperature) and the measured physical reality. The engine mints Value Tokens to human citizens or actuates automated relays to apply the control variable $u(t)$ (physical labor or electrical power) to correct the error.

\section{Logistical Aggregation and Network Physics}
Hard representations are not constrained to localized thermodynamics; they encompass macroscopic logistics and the flow of physical goods through the mesh network. In scenarios where sovereign fabrication is unfeasible (e.g., printed circuit boards, heavy metals), the system must interface with the legacy environment to procure materials (Scenario Alpha: Fabrication).

\subsection{Economic Order Quantity and Delay Capacitors}
The system models procurement mathematically to minimize the thermodynamic and exergy costs of shipping and packaging (the "Ecological Leeching" penalty). This is achieved through decentralized batching algorithms derived from the Economic Order Quantity (EOQ) model:
\begin{equation}
EOQ = \sqrt{\frac{2DS}{H}}
\end{equation}
where $D$ is the aggregated demand rate across the node, $S$ is the fixed cost per order (which in Oasis includes the exact exergy cost of freight transport, not just fiat), and $H$ is the thermodynamic holding or warehousing cost.

By pooling individual citizen requests via the orchestrator and executing bulk purchases, the simulation forces temporal delays. The time-delay acts as a \textit{virtual capacitor}, storing demand potential over days or weeks until the thermodynamic efficiency criteria for a physical order are met. This explicitly prevents the high-entropy, instant-gratification delivery mechanics prevalent in the legacy fiat economy \cite{harris1913how}.

\vspace{1cm}
\begin{thebibliography}{99}
\bibitem{shannon1949communication}
Shannon, C. E. (1949). \textit{Communication in the presence of noise}. Proceedings of the IRE, 37(1), 10-21.

\bibitem{harris1913how}
Harris, F. W. (1913). \textit{How many parts to make at once}. Factory, The Magazine of Management, 10(2), 135-136.
\end{thebibliography}
